% GENERATED FILE. Edit canonical lessons, then run npm run generate:books.
% canonical-source-hash: fnv1a64:7220598c0456d683
% canonical-lessons: HI-C409-eyarport, HI-C409-viman, HI-C409-payalat, HI-C409-uran, HI-C409-agaman

\chapter{Airports and Flights}
\label{ch:hi-airports-and-flights}
\hlchaptermodality{409}

\begin{chapteropening}
\textbf{By the end of this chapter:} \emph{I can talk about the airport, the plane and the pilot, and read the arrivals board.}

Complete the last lesson of chapter 409: I can talk about the airport, the plane and the pilot, and read the arrivals board.
\end{chapteropening}

\section[\texorpdfstring{\hi{एयरपोर्ट} (eyarporṭ)}{एयरपोर्ट (eyarporṭ)}]{\hi{एयरपोर्ट} (eyarporṭ) — an airport}

\label{lesson:HI-C409-eyarport}



\begin{quote}
Before the new one: say the Hindi for a path, then the Hindi for to chase off.
\end{quote}

\subsection*{You'll want to know: \hi{एयरपोर्ट}}
\textbf{\hi{एयरपोर्ट}} — \emph{eyarporṭ} — \textquotedblleft{}an airport\textquotedblright{}.

\begin{quote}
\hi{एयरपोर्ट} \hi{शहर} \hi{से} \hi{काफ़ी} \hi{दूर} \hi{है।}
\end{quote}

\begin{quote}
\emph{eyarporṭ śahar se kāfī dūr hai.}
\end{quote}

\begin{quote}
The airport is quite far from the city.
\end{quote}

Masculine.

\begin{cousinweb}
\textbf{\hi{एयरपोर्ट}} (\emph{eyarporṭ}) is borrowed from English \emph{airport}, the word most people use in speech.
\end{cousinweb}

\begin{grammarlens}[title={What you've built}]
One more word for getting around.
\end{grammarlens}

\subsection*{Your turn}
\begin{itemize}
\raggedright
  \item \emph{Say it:} \emph{eyarporṭ}
  \item \emph{Say it:} \emph{eyarporṭ}, once more
  \item \emph{Say it:} \emph{eyarporṭ śahar se kāfī dūr hai.} — The airport is quite far from the city.
  \item \emph{Recall:} say the Hindi for a path, then the Hindi for to chase off, then say \emph{eyarporṭ} again
\end{itemize}

\begin{tcolorbox}[breakable,colback=teal!4,colframe=teal!35!black,title={Before you move on}]
What does \hi{एयरपोर्ट} mean? (\textquotedblleft{}An airport\textquotedblright{}.) Say it once more.
\end{tcolorbox}
\section[\texorpdfstring{\hi{विमान} (vimān)}{विमान (vimān)}]{\hi{विमान} (vimān) — an aircraft (formal)}

\label{lesson:HI-C409-viman}



\begin{quote}
Before the new one: say the Hindi for to chase off, then the Hindi for an airport.
\end{quote}

\subsection*{You'll want to know: \hi{विमान}}
\textbf{\hi{विमान}} — \emph{vimān} — \textquotedblleft{}an aircraft (formal)\textquotedblright{}.

\begin{quote}
\hi{विमान} \hi{बादलों} \hi{के} \hi{ऊपर} \hi{उड़} \hi{रहा} \hi{था।}
\end{quote}

\begin{quote}
\emph{vimān bādloṁ ke ūpar uṛ rahā thā.}
\end{quote}

\begin{quote}
The plane was flying above the clouds.
\end{quote}

Masculine.

\begin{cousinweb}
\textbf{\hi{विमान}} (\emph{vimān}) is a tatsama, Sanskrit \emph{vimāna}, the self-moving chariot of the gods in the old stories and later any vehicle (Monier-Williams); modern Hindi gave it the aircraft.
\end{cousinweb}

\begin{grammarlens}[title={What you've built}]
One more word for getting around.
\end{grammarlens}

\subsection*{Your turn}
\begin{itemize}
\raggedright
  \item \emph{Say it:} \emph{vimān}
  \item \emph{Say it:} \emph{vimān}, once more
  \item \emph{Say it:} \emph{vimān bādloṁ ke ūpar uṛ rahā thā.} — The plane was flying above the clouds.
  \item \emph{Recall:} say the Hindi for to chase off, then the Hindi for an airport, then say \emph{vimān} again
\end{itemize}

\begin{tcolorbox}[breakable,colback=teal!4,colframe=teal!35!black,title={Before you move on}]
What does \hi{विमान} mean? (\textquotedblleft{}An aircraft (formal)\textquotedblright{}.) Say it once more.
\end{tcolorbox}
\section[\texorpdfstring{\hi{पायलट} (pāyalaṭ)}{पायलट (pāyalaṭ)}]{\hi{पायलट} (pāyalaṭ) — a pilot}

\label{lesson:HI-C409-payalat}



\begin{quote}
Before the new one: say the Hindi for an airport, then the Hindi for an aircraft.
\end{quote}

\subsection*{You'll want to know: \hi{पायलट}}
\textbf{\hi{पायलट}} — \emph{pāyalaṭ} — \textquotedblleft{}a pilot\textquotedblright{}.

\begin{quote}
\hi{पायलट} \hi{ने} \hi{कहा} \hi{कि} \hi{मौसम} \hi{ठीक} \hi{है।}
\end{quote}

\begin{quote}
\emph{pāyalaṭ ne kahā ki mausam ṭhīk hai.}
\end{quote}

\begin{quote}
The pilot said that the weather was fine.
\end{quote}

Masculine or feminine, by the person meant.

\begin{cousinweb}
\textbf{\hi{पायलट}} (\emph{pāyalaṭ}) is borrowed from English \emph{pilot}.
\end{cousinweb}

\begin{grammarlens}[title={What you've built}]
One more word for getting around.
\end{grammarlens}

\subsection*{Your turn}
\begin{itemize}
\raggedright
  \item \emph{Say it:} \emph{pāyalaṭ}
  \item \emph{Say it:} \emph{pāyalaṭ}, once more
  \item \emph{Say it:} \emph{pāyalaṭ ne kahā ki mausam ṭhīk hai.} — The pilot said that the weather was fine.
  \item \emph{Recall:} say the Hindi for an airport, then the Hindi for an aircraft, then say \emph{pāyalaṭ} again
\end{itemize}

\begin{tcolorbox}[breakable,colback=teal!4,colframe=teal!35!black,title={Before you move on}]
What does \hi{पायलट} mean? (\textquotedblleft{}A pilot\textquotedblright{}.) Say it once more.
\end{tcolorbox}
\section[\texorpdfstring{\hi{उड़ान} (uṛān)}{उड़ान (uṛān)}]{\hi{उड़ान} (uṛān) — a flight}

\label{lesson:HI-C409-uran}



\begin{quote}
Before the new one: say the Hindi for an aircraft, then the Hindi for a pilot.
\end{quote}

\subsection*{You'll want to know: \hi{उड़ान}}
\textbf{\hi{उड़ान}} — \emph{uṛān} — \textquotedblleft{}a flight\textquotedblright{}.

\begin{quote}
\hi{मेरी} \hi{उड़ान} \hi{रात} \hi{दस} \hi{बजे} \hi{है।}
\end{quote}

\begin{quote}
\emph{merī uṛān rāt das baje hai.}
\end{quote}

\begin{quote}
My flight is at ten at night.
\end{quote}

Feminine.

\begin{cousinweb}
\textbf{\hi{उड़ान}} (\emph{uṛān}) is the noun of \hi{उड़ना}, \textquotedblleft{}to fly\textquotedblright{}, which goes back to Sanskrit \emph{uḍḍīyate}, \textquotedblleft{}flies up\textquotedblright{} (CDIAL 1697); CDIAL does not list the noun itself.
\end{cousinweb}

\begin{grammarlens}[title={What you've built}]
One more word for getting around.
\end{grammarlens}

\subsection*{Your turn}
\begin{itemize}
\raggedright
  \item \emph{Say it:} \emph{uṛān}
  \item \emph{Say it:} \emph{uṛān}, once more
  \item \emph{Say it:} \emph{merī uṛān rāt das baje hai.} — My flight is at ten at night.
  \item \emph{Recall:} say the Hindi for an aircraft, then the Hindi for a pilot, then say \emph{uṛān} again
\end{itemize}

\begin{tcolorbox}[breakable,colback=teal!4,colframe=teal!35!black,title={Before you move on}]
What does \hi{उड़ान} mean? (\textquotedblleft{}A flight\textquotedblright{}.) Say it once more.
\end{tcolorbox}
\section[\texorpdfstring{\hi{आगमन} (āgaman)}{आगमन (āgaman)}]{\hi{आगमन} (āgaman) — arrival (formal)}

\label{lesson:HI-C409-agaman}



\begin{quote}
Before the new one: say the Hindi for a pilot, then the Hindi for a flight.
\end{quote}

\subsection*{You'll want to know: \hi{आगमन}}
\textbf{\hi{आगमन}} — \emph{āgaman} — \textquotedblleft{}arrival (formal)\textquotedblright{}.

\begin{quote}
\hi{विमान} \hi{का} \hi{आगमन} \hi{सुबह} \hi{सात} \hi{बजे} \hi{है।}
\end{quote}

\begin{quote}
\emph{vimān kā āgaman subah sāt baje hai.}
\end{quote}

\begin{quote}
The plane arrives at seven in the morning.
\end{quote}

Masculine.

\begin{cousinweb}
\textbf{\hi{आगमन}} (\emph{āgaman}) is a tatsama, Sanskrit \emph{āgamana}, \textquotedblleft{}coming, arriving\textquotedblright{} (Monier-Williams), from the root \emph{gam}, \textquotedblleft{}go\textquotedblright{}: the word on station and airport boards.
\end{cousinweb}

\begin{grammarlens}[title={What you've built}]
One more word for getting around.
\end{grammarlens}

\subsection*{Your turn}
\begin{itemize}
\raggedright
  \item \emph{Say it:} \emph{āgaman}
  \item \emph{Say it:} \emph{āgaman}, once more
  \item \emph{Say it:} \emph{vimān kā āgaman subah sāt baje hai.} — The plane arrives at seven in the morning.
  \item \emph{Recall:} say the Hindi for a pilot, then the Hindi for a flight, then say \emph{āgaman} again
\end{itemize}

\begin{tcolorbox}[breakable,colback=teal!4,colframe=teal!35!black,title={Before you move on}]
What does \hi{आगमन} mean? (\textquotedblleft{}Arrival (formal)\textquotedblright{}.) Say it once more.
\end{tcolorbox}
